\chapter{Vecteurs, bases et repères de l’espace}
\section*{Exercice 5.1 -- Cube}
$E=(0,0,1)$, $C=(1,1,0)$, donc $\overrightarrow{EC}=(1,1,-1)$. Avec $B=(1,0,0)$ et $H=(0,1,1)$, le milieu est $(1/2,1/2,1/2)$, centre du cube.

\section*{Exercice 5.2 -- Combinaison}
$au+bv+cw=(a+b,b,c)$, donc $b=1$, $a=2$, $c=2$. La décomposition est $2u+v+2w$ ; on vérifie en additionnant les coordonnées.

\section*{Exercice 5.3 -- Plan ou droite}
Non : si $v=2u$ et $u\ne0$, les points sont $A+(s+2t)u$, donc forment une droite. Si les deux directions sont nulles, il ne reste qu'un point. La non-colinéarité est essentielle.

\section*{Exercice 5.4 -- Coplanarité}
Le troisième vaut deux fois le premier plus trois fois le deuxième. Ils sont donc coplanaires et ne constituent pas une base de l'espace.

\section*{Exercice 5.5 -- Droites gauches}
Elles sont disjointes puisque les coordonnées $z$ valent $0$ et $1$ ; leurs directions $(1,0,0)$ et $(0,1,0)$ ne sont pas colinéaires. Deux droites coplanaires non parallèles se rencontreraient ; elles sont donc gauches.

\section*{Exercice 5.6 -- Unicité}
Cette relation donnerait deux écritures du vecteur nul : coefficients $(0,0,0)$ et $(1,2,-3)$. L'unicité des coordonnées dans une base les rend incompatibles. La relation exprimerait aussi $w=(u+2v)/3$, donc la coplanarité.
